% ASim 2026 full paper
% Main manuscript: limited to eight pages including references.
% Material held for a possible supplement is stored in appendix-draft.tex.

\documentclass[12pt,a4paper]{article}
\usepackage{Asim2026}
\usepackage{microtype}

\addbibresource{references.bib}
\graphicspath{{figures/}}

\setcounter{secnumdepth}{0}
\setlength{\textfloatsep}{10pt plus 2pt minus 2pt}
\setlength{\floatsep}{10pt plus 2pt minus 2pt}
\setlength{\intextsep}{10pt plus 2pt minus 2pt}
\renewcommand{\topfraction}{0.95}
\renewcommand{\bottomfraction}{0.90}
\renewcommand{\textfraction}{0.05}
\renewcommand{\floatpagefraction}{0.85}

\title{Comparative Assessment of Group Comfort Model Resolution in HVAC Control Using High-Resolution Occupancy Tracking}
\date{}

\begin{document}
\maketitle

\noindent
K.~Horikoshi$^{1}$,
E.~Ono$^{2}$,
K.~Mihara$^{3}$,
A.~Chong$^{1,*}$

\renewcommand{\thefootnote}{\fnsymbol{footnote}}
\footnotetext[1]{Corresponding author email: adrian.chong@nus.edu.sg}
\renewcommand{\thefootnote}{\arabic{footnote}}

\vspace{\baselineskip}

\noindent
$^{1}$Department of the Built Environment, National University of Singapore, Singapore\\
$^{2}$Kajima Technical Research Institute Singapore, Kajima Corporation, Singapore\\
$^{3}$Kajima Technical Research Institute, Kajima Corporation, Tokyo, Japan

\noindent\textbf{ABSTRACT}
Occupant-centric HVAC control in shared offices must aggregate several personal comfort models into one room-level setpoint, yet hybrid and activity-based work change both the number and composition of occupants present.
This study quantifies when daily or real-time group comfort models justify their additional occupancy information.
Personal comfort models were developed for 39 participants from over 2,000 thermal-preference responses collected over four weeks in a Singapore office building, and an occupancy-log-based Monte Carlo simulation compared a static group model (regular room members), a daily model (daily attendees), a real-time model (occupants present at each interval), and a fixed 24$^\circ$C baseline.
Across four rooms, the real-time model achieved the highest mean comfort probability at every tested group size and, at the largest sizes,
exceeded the baseline by 12--15 and daily aggregation by 6--8 percentage
points; the static model approached the baseline as group size increased.
Real-time setpoint changes became smaller as mean occupancy increased but remained about 0.8$^\circ$C at the upper end of the observed daily mean-occupancy range, and the share of control intervals requiring an actionable setpoint step rose from about 37\% at low occupant turnover to roughly 90\% near complete turnover.
These findings demonstrate that group size, daily mean occupancy, and occupant turnover jointly shape the practical value of higher-resolution comfort aggregation, providing criteria for selecting group comfort models in shared offices.

\vspace{\baselineskip}
\noindent\textbf{KEYWORDS}
\noindent
Occupant-centric control, Occupancy dynamics, Thermal comfort, Agent simulation

\section{INTRODUCTION}

Occupant-centric control seeks to operate building systems using information about the people who occupy them, rather than applying one fixed condition to an assumed average occupant.
Real-world implementation, however, remains limited by data availability, scalability, and operational complexity \parencite{xie_review_2020,soleimanijavid_challenges_2024}.
These barriers become especially important in shared offices, where individual comfort information must ultimately be translated into a single heating, ventilation, and air-conditioning (HVAC) setpoint for several people.
A personal comfort model (PCM) predicts an individual's thermal response.
In a multi-occupant zone, PCMs must be combined into a group comfort model (GCM).
The resulting setpoint depends on both the PCM formulation and on who is represented when the decision is made.
Previous multi-occupant simulations showed that the comfort-side benefit of integrating personal comfort information diminishes as more occupants share a thermal zone, because heterogeneous preferences are averaged into a group representation \parencite{jung_energy_2020, wang_enhancing_2026}.
% At larger group scales, recent stochastic simulations similarly showed that increasing group size and internal preference diversity reduced the potential maximum satisfaction rate, while larger groups exhibited lower thermal sensitivity \parencite{wang_enhancing_2026}.
The usefulness of a high-resolution comfort model also depends on whether the HVAC system can act at a compatible spatial and control resolution \parencite{ono_effects_2022}.
This aggregation problem is particularly important under dynamic occupancy, where arrivals, departures, and movement change both occupant count and composition \parencite{marzban_review_2023}.
A static GCM representing all regular room members may therefore be poorly matched to sparse attendance, whereas a real-time GCM can update the represented comfort distribution whenever the occupant composition changes.
Real-time occupancy tracking is becoming increasingly available in office practice through IoT and indoor-positioning systems \parencite{zafari_survey_2019}.
However, using these data for real-time comfort aggregation adds sensing, computation, and control-update requirements.
Its value should therefore be judged against simpler alternatives, not merely against a fixed thermostat.
This paper compares three temporal resolutions of comfort aggregation: a static GCM (SGCM) based on regular room members, a daily GCM (DGCM) based on attendees on a given day, and a real-time GCM (RT-GCM) based on the occupants present at each interval.
Using observed office movement records and field-derived PCMs, the study addresses three questions: (1) how the comfort benefit of each GCM changes with the sampled group size; (2) how realized occupancy affects the magnitude of real-time setpoint adjustments; and (3) how occupant-composition turnover affects the frequency of actionable setpoint updates.

\section{METHOD}
Figure~\ref{fig:simulation_flow} summarizes the study workflow.
Field records provide two complementary inputs: thermal-preference votes and indoor conditions for fitting PCMs, and room-level entry/exit records that preserve observed attendance and movement as location-log agents.
For each of four target rooms and each sampled subgroup size, PCMs are randomly paired with location-log agents, whose recorded four-week stays are replayed in Monte Carlo trials.
Static, daily, and real-time group comfort models then select setpoints from the same PCM pool at their respective aggregation resolutions.
Person-time comfort probability and the control consequences of real-time updating are evaluated across rooms and sampled group sizes.

\begin{figure}[!htbp]
  \centering
  \includegraphics[width=\linewidth]{figures/SimulationFlow_small.png}
  % \includegraphics[width=0.67\linewidth]{figures/SimulationFlow_wider.png}
  \caption{Workflow for the occupancy-log-based Monte Carlo comparison of static, daily, and real-time group comfort models.}
  \label{fig:simulation_flow}
\end{figure}

\subsection{Field data}
Data were collected from 29 September to 24 October 2025 in an office building in Singapore.
The observed rooms were served by a central HVAC system with zone-level VAV control and summarized in~\autoref{table:RoomList}.
\begin{table}[htbp!]
  \centering
  \small
  \caption{Simulation target rooms}
  \begin{tabular}{@{}
    p{2.6cm}
    p{3.8cm}
    M{2.2cm}
    M{1.8cm}
    M{1.8cm}
    @{}}
    \toprule
    \TableHeader{Room} &
    \TableHeaderBreak{Occupants' \TableHeaderLineBreak Job function} &
    \TableHeaderBreak{Seating type} &
    \TableHeaderBreak{Floor area\TableHeaderLineBreak \((\mathrm{m^2})\)} &
    \TableHeaderBreak{Seating\TableHeaderLineBreak capacity} \\
    \midrule
    Office 5A & Real estate development & Assigned & 221.40 & 15 \\
    Office 5B & Research \& development & Assigned & 163.56 & 18 \\
    Shared Office 5A & Real estate development & Shared & 221.40 & 14 \\
    Shared Office 5B & Building construction   & Shared & 213.73 & 20 \\
    \bottomrule
\end{tabular}
  \label{table:RoomList}
\end{table}
Data from 39 participants with sufficiently complete survey records, together with location records for regular room users, were retained.
The protocol was approved by the National University of Singapore Institutional Review Board (NUS-IRB-2025-498) and informed consent was obtained.
Thermal-preference votes were collected by web survey up to six times per workday using the three ASHRAE Thermal Preference categories ``warmer'', ``no change'', and ``cooler'' \parencite{ashrae_ansiashrae_2022}.
Responses within 20 minutes of room entry were excluded to reduce transient exposure effects.
In total, 2,608 responses were available.
Room air temperature and relative humidity were measured while setpoints were varied between 22$^\circ$C and 27$^\circ$C to expose participants to an informative range of conditions.
Anonymized location records were obtained from a face-recognition system installed at room entrances and in corridors and reconstructed into continuous room stays with entry and exit times.

\subsection{Personal and group comfort models}
An ordered multinomial logistic PCM was fitted for each participant using room air temperature as the predictor and thermal preference as the response, following a prior agent-based control study \parencite{jung_energy_2020}.
Let $p_i(T)$ be the PCM probability that occupant $i$ votes ``No change'' at temperature $T$.
For any represented group $G$, the group comfort curve and selected setpoint were

\begin{equation}
 \bar p_G(T)=\frac{1}{|G|}\sum_{i\in G}p_i(T),
 \qquad
 T_G^*=\operatorname*{arg\,max}_{T\in\mathcal{T}_{\mathrm{set}}}\bar p_G(T),
 \label{eq:gcm}
\end{equation}

where $\mathcal{T}_{\mathrm{set}}$ is the candidate setpoint grid (0.1$^\circ$C resolution).
All GCM policies used the same optimization rule; only the group represented by $G$ changed (Table~\ref{tab:policies}).
% This isolates the value of attendance resolution from the setpoint-selection algorithm.

\begin{table}[!htbp]
  \caption{Compared comfort-representation policies.}
  \label{tab:policies}
  \begin{tabularx}{\linewidth}{>{\raggedright\arraybackslash}p{0.19\linewidth}
                                >{\raggedright\arraybackslash}p{0.27\linewidth}
                                >{\raggedright\arraybackslash}X}
    \toprule
    \textit{Policy} & \textit{Represented group} & \textit{Update interpretation} \\
    \midrule
    Static GCM (SGCM) & Sampled regular room members & One group curve for the study period \\
    Daily GCM (DGCM) & Sampled members attending that day & Attendance-aware curve updated daily \\
    Real-time GCM (RT-GCM) & Sampled members present at the interval & Composition-aware curve updated during the day \\
    Baseline & None & Fixed 24$^\circ$C setpoint \\
    \bottomrule
  \end{tabularx}
\end{table}

\subsection{Occupancy-log-based simulation and metrics}
Each trial sampled $K$ location-log agents from a room's regular-member
pool and randomly paired them with $K$ PCMs, randomly drawn without replacement from the pool of 39 PCMs; $K$ denotes the sampled
group size.
The agents retained their observed stay logs in their original rooms; random PCM assignment varied the comfort-profile composition without inventing movement trajectories.
The present set $S_{r,t}$ and instantaneous occupancy $N_{r,t}=|S_{r,t}|$ varied over time.
Decision-time occupancy density was $N_{r,t}/K$; daily mean occupancy $\bar N_d$ was the mean of $N_{r,t}$ over occupied intervals on day $d$, and $\bar N_d/K$ its daily normalized counterpart.
For each room and each feasible $K$, 1,000 trials were simulated over the four-week record.
At every occupied interval, each policy's setpoint was scored against the PCMs of the people actually present.
Mean comfort probability was defined as

\begin{equation}
 \mathrm{CP}_{r,d}^{(\pi,q)}
 =\frac{\displaystyle\sum_{t\in\mathcal{T}_{r,d}}
       \sum_{i\in S_{r,t}^{(q)}}p_i\!\left(T_{r,t}^{(\pi,q)}\right)}
      {\displaystyle\sum_{t\in\mathcal{T}_{r,d}}\left|S_{r,t}^{(q)}\right|}
 % \qquad
 % \overline{\mathrm{CP}}_{r}^{(\pi)}
 % =\frac{1}{Q}\sum_{q=1}^{Q}
 %  \frac{1}{\left|\mathcal{D}_{r}^{(q)}\right|}
 %  \sum_{d\in\mathcal{D}_{r}^{(q)}}\mathrm{CP}_{r,d}^{(\pi,q)}.
 \label{eq:comfort_metric}
\end{equation}

Here, $\pi$ denotes the policy, $q$ the Monte Carlo trial, $\mathcal{T}_{r,d}$ the occupied control intervals for room $r$ on day $d$.
% $\mathcal{D}_{r}^{(q)}$ the valid room-days.
% and $Q=1{,}000$
Room-day values were averaged across days within each trial and then summarized across the 1,000 Monte Carlo trials.
Improvement is reported as a percentage-point difference from the fixed 24$^\circ$C baseline.
Control consequences were summarized by the absolute RT-GCM setpoint change between consecutive intervals and by an actionable-change indicator.
Occupant-composition change was measured using Jaccard turnover.
For a lag $\Delta t$, the occupant turnover rate (UTR) was

\begin{equation}
 \mathrm{UTR}_{r,t}=1-\frac{|S_{r,t}\cap S_{r,t-\Delta t}|}
 {|S_{r,t}\cup S_{r,t-\Delta t}|}.
 \label{eq:utr}
\end{equation}

UTR is zero when the two occupant sets are identical and one when they have no members in common.
The main action-frequency result uses 15-minute transitions.
% A transition was classified as actionable when the setpoint changed by at least 0.5$^\circ$C.

\section{RESULTS}
This section compares comfort performance and time-varying occupancy density (Figures~\ref{fig:comfort} and~\ref{fig:OccuDensity}), then examines the magnitude and frequency of real-time setpoint changes (Figures~\ref{fig:control} and~\ref{fig:turnover_action}).

\subsection{Comfort performance across aggregation resolutions}
\begin{figure}[!htbp]
  \centering
  \includegraphics[width=0.98\linewidth]{fig1_analysis23nr_with_fixed24_raw2000_64_100_all_solid_small_markers_paper}
  % {fig1_comfort_probability_64_100_paper}
  \caption{Mean person-time comfort probability under the three GCM policies and fixed baseline.}
  \label{fig:comfort}
\end{figure}
\begin{figure}[!htbp]
  \centering
  \includegraphics[width=0.98\linewidth]{occupancy_density_timeofday_k12_paper}
  \caption{Decision-time occupancy density $N_{r,t}/K$ by time of day; lines show means, with 95\% confidence and simulation intervals}
  \label{fig:OccuDensity}
\end{figure}
Figure~\ref{fig:comfort} shows the same policy ranking in all four analyzed rooms: RT-GCM led at every $K$, followed by DGCM, while SGCM tracked near the fixed 24$^\circ$C baseline, steady at 65--70\% regardless of $K$.
RT-GCM declined from 96--97\% at $K=2$ to 79--82\% at the largest tested $K$; DGCM declined from 91--94\% to 72--76\%; SGCM fell from 72--73\% to 64--66\%, converging on the baseline.
Figure~\ref{fig:OccuDensity} shows that mean decision-time occupancy density remained below full occupancy and varied by room and time.
Thus, large member pools often produced partial attendance, allowing the present group to differ from regular membership and daily attendees and preserving an advantage for attendance-aware aggregation.

\subsection{Setpoint adjustments and occupant turnover}

Figure~\ref{fig:control}(a) relates daily mean occupancy to the magnitude of RT-GCM setpoint changes when a change occurred.
The central trend decreased from about 1.3$^\circ$C near a daily mean occupancy of one to roughly 1.0$^\circ$C near four and 0.8$^\circ$C at the highest observed daily mean occupancy, with the widest bands at low occupancy.
The typical step remained above the common 0.5$^\circ$C thermostat resolution.
Figure~\ref{fig:control}(b) resolves the same changed-step magnitude by $K$ and normalized daily mean occupancy $\bar N_d/K$.
The larger steps were concentrated at low normalized occupancy and generally declined as occupancy increased across $K$, consistent with Figure~\ref{fig:control}(a).

\begin{figure}[!htbp]
  \centering
  \begin{minipage}[t]{0.49\linewidth}
    \centering
    \includegraphics[width=\linewidth]{figures/fig3a_setpoint_adjustment_by_occupancy_compact_paper.png}
    \parbox{0.98\linewidth}{\centering (a) Mean absolute setpoint adjustment versus daily mean occupancy}
  \end{minipage}\hfill
  \begin{minipage}[t]{0.49\linewidth}
    \centering
    \includegraphics[width=\linewidth]{figures/fig3b_setpoint_adjustment_by_group_size_and_occupancy_density_compact_paper.png}
    \parbox{0.98\linewidth}{\centering (b) Mean absolute setpoint adjustment across $K$ and normalized daily mean occupancy, $\bar{N}_d/K$}
  \end{minipage}
    \caption{RT-GCM setpoint-change magnitude across occupancy conditions. (a) Daily mean absolute  setpoint change, versus daily mean occupancy; bands denote the 25--75\% and 10--90\% percentile ranges. (b) The same metric across $K$ and normalized daily mean occupancy $\bar N_d/K$; blank cells indicate insufficient observations.}
  \label{fig:control}
\end{figure}

Figure~\ref{fig:turnover_action}(a) relates occupant turnover to the frequency of actionable RT-GCM updates (at least 0.5$^\circ$C).
The horizontal axis is the 15-minute UTR, where zero indicates no change in occupant composition and one indicates complete replacement.
The vertical axis is the percentage of consecutive control intervals in which the RT-GCM setpoint changed by at least 0.5$^\circ$C.
The actionable-step rate increased monotonically from about 37\% at UTR near 0.05 to 63--69\% at 0.4, 77--78\% at 0.65, and 89--93\% at 0.9--1.0.
The near-overlap of the $K$ bands indicates that occupant replacement, rather than group size, was the dominant driver of short-interval action frequency.
Shared offices turned over dynamically through office hours, while the assigned-seat office (5A) stayed near 0.2--0.4 UTR and rarely reached the highest bins.
Figure~\ref{fig:turnover_action}(b) resolves the temporal patterns for two large-$K$ room examples at 30-minute resolution.
Office 5A showed pronounced Monday-afternoon turnover, reaching a UTR of 1.0, whereas Shared Office 5B exhibited intermittent peaks across the week, including a Monday-afternoon peak near 0.8.
These profiles show that high turnover was concentrated in particular room-time combinations rather than being a uniform property of all occupied intervals.
% Estimates in the highest-UTR bins rest on comparatively few transitions.

\begin{figure}[!htbp]
  \centering
  \begin{minipage}[t]{0.49\linewidth}
    \centering
    \includegraphics[width=\linewidth]{figures/fig4a_action_rate_by_turnover_compact_paper.png}
    \parbox[t]{0.96\linewidth}{\centering (a) Actionable RT-GCM update rate versus 15-minute UTR.}
  \end{minipage}\hfill
  \begin{minipage}[t]{0.49\linewidth}
    \centering
    \includegraphics[width=\linewidth]{figures/fig4b_weekly_30min_utr_selected_compact_paper.png}
    \parbox[t]{0.96\linewidth}{\centering (b) Weekday 15-minute UTR profiles for selected room examples.}
  \end{minipage}
  \caption{Occupant turnover and actionable RT-GCM updates. (a) Actionable-update rate versus 15-minute UTR; an update is actionable when $|\Delta T|\geq0.5\,^{\circ}\mathrm{C}$. Lines denote bins of $K$. (b) Weekday 15-minute UTR profiles for Office 5A ($K=12$) and Shared Office 5B ($K=8$).}
  \label{fig:turnover_action}
\end{figure}


\section{DISCUSSION}
Using occupancy logs and field-derived PCMs under a shared setpoint-selection rule, this study shows that the value of temporal resolution depends on $K$, daily mean occupancy, and UTR.
These metrics can identify where real-time tracking may justify its cost.

\subsection{Why aggregation resolution changes the group-size result}
Prior simulations showed that the benefits of PCM integration diminish as more occupants share a zone, particularly beyond six occupants \parencite{jung_energy_2020}; at larger group scales, increasing group size and internal preference diversity likewise reduced the potential maximum satisfaction rate, while larger groups became less thermally sensitive \parencite{wang_enhancing_2026}.
The SGCM here follows the same directional pattern, crossing zero at $K\approx6$--8.
DGCM and RT-GCM, however, retained positive improvement at the same $K$ by excluding absent members.
This study therefore extends prior group-size findings by showing that comfort-aggregation outcomes depend not only on potential group size but also on realized occupancy (Figure~\ref{fig:OccuDensity}) and dynamic changes in occupant composition.
RT-GCM also exceeded DGCM by 6--8 percentage points after the initial rise with $K$, reflecting the value of knowing who is present now rather than who attended that day.
These predicted gains assume immediate setpoint selection and remain conditional on sensing reliability, HVAC response, and occupant adaptation.
% This is not proof that every real-time system will deliver the same field improvement: the metric is a predicted mean ``No change'' probability under a shared setpoint with immediate setpoint selection, so it identifies a theoretical value of temporal resolution whose actual delivery depends on sensing reliability, response, and adaptation.

\subsection{From comfort information to actionable control}
A high-resolution comfort representation is useful only when the HVAC system can respond to the resulting setpoint differences, reflecting the mismatch between comfort-model and control resolution reported by \textcite{ono_effects_2022}.
Although the simulation used a 0.1$^\circ$C grid, typical changed steps remained near 0.8$^\circ$C at higher mean occupancy and exceeded the 0.5$^\circ$C actionable threshold.
This indicates actuation-relevant potential rather than delivered field benefit; reviews report limited field implementation and persistent data, computation, BAS-integration, and scalability barriers \parencite{xie_review_2020,soleimanijavid_challenges_2024}.
UTR complements mean occupancy by indicating whether represented comfort profiles are being replaced, extending prior discussions of dynamic diversity in comfort-driven control \parencite{jung_energy_2020}.
The near-overlap across $K$ in Figure~\ref{fig:turnover_action}(a) shows that turnover is a better guide to update frequency than $K$ alone.
The three aggregation levels can consequently be treated as a deployment hierarchy: SGCM is defensible in stable assigned-seat zones, where UTR remained near 0.2--0.4 through office hours and the regular-member group matched presence; DGCM is a lower-complexity compromise when daily attendance varies but within-day composition stays relatively stable.
RT-GCM is most valuable where UTR instead fluctuated dynamically, with sparse attendance and frequent movement.
UTR can gate real-time updates to high-turnover periods, with deadbands or minimum hold times limiting excessive cycling.
Such gating does not remove sensing trade-offs in accuracy, latency, energy use, cost, scalability, and privacy \parencite{zafari_survey_2019}.
Because activity-based working outcomes depend on implementation and occupant use \parencite{marzban_review_2023}, these policy thresholds should be calibrated to each workplace.
% Even there, UTR can gate real-time updates to high-turnover periods, with deadbands or minimum hold times preventing excessive cycling.

\subsection{Limitations and research needs}
The results are conditional on one Singapore office building, four weeks, four rooms, and 39 PCMs; random PCM assignment broadened the combinations, but comfort remained predicted rather than closed-loop.
Energy use, HVAC dynamics and integration cost were not evaluated, so the deployment hierarchy should be interpreted as a comfort-based screening guide.
Mean comfort probability also cannot show whether lower aggregate comfort reflects minority disadvantage or a low attainable maximum under preference diversity; future work should evaluate occupant-level fairness and group-relative performance \parencite{wang_enhancing_2026}.


\section{CONCLUSION AND IMPLICATIONS}
This study compared static, daily, and real-time comfort aggregation under observed office occupancy dynamics.
RT-GCM achieved the highest predicted mean comfort probability across all tested rooms and group sizes.
SGCM improvement diminished to approximately zero as the group size grew, whereas DGCM and RT-GCM remained beneficial by excluding absent occupants.
The additional RT-GCM gain over DGCM was 6--8 percentage points, showing that within-day composition information can matter even after daily attendance is known.
Real-time information also implied operational consequences.
Setpoint changes shrank as daily mean occupancy increased but still exceeded a common 0.5$^\circ$C control step.
The actionable-update rate rose sharply with UTR but was insensitive to $K$.
Within this comfort-only simulation scope, the practical implication is selective deployment: static aggregation for stable membership, daily aggregation for variable attendance with moderate turnover, and real-time tracking where sparse, rapidly changing composition justifies the sensing and control effort.
UTR could serve as a screening metric for selecting the aggregation resolution.

\section*{ACKNOWLEDGEMENTS}

This research is supported by Kajima Corporation (grant number A-8001911-00-00).  

\section*{REFERENCES}

\printbibliography[heading=none]

\end{document}
